\chapter{On-Chain Order Books and Perpetual DEXs}\label{ch:m3:on-chain-order-books-and-perpetual-dexs}

On one decentralised venue every order, every cancel and every liquidation of a \omterm{def:m3:perpetual-futures:perp}{perpetual futures} market
is a transaction in a block. The venue's \omterm{def:m3:blockchains-for-traders:pow}{validators} agree not only on what happened but on the order in
which it happened, and that order decides who is filled: a market maker whose quote has gone stale sends a
cancel, an arbitrageur sends an order to hit the quote, and both arrive in the same block. On a
\omterm{def:m3:centralised-exchanges:cex}{centralised exchange} the matching engine's clock settles the race. On a chain, a rule written into the
protocol does. This chapter explains why putting an order book on a chain is hard, examines in detail the
design of the venue that has gone furthest in doing it, looks at what went wrong when markets were
stressed, and compares the other designs: order books kept off chain by \omterm{def:m3:blockchains-for-traders:pow}{validators}, and exchanges that
dispense with order books altogether and trade at an \omterm{def:m3:blockchains-for-traders:contract}{oracle}'s price.

\section{Why an on-chain order book is hard}

\begin{definition}[Perpetual DEX, on-chain order book]\label{def:m3:on-chain-order-books-and-perpetual-dexs:perpdex}
A \emph{perpetual DEX}\index{perpetual DEX} is a \omterm{def:m3:automated-market-makers:amm}{decentralised exchange} for \omterm{def:m3:perpetual-futures:perp}{perpetual futures}, whose margin,
positions and settlement are held and computed by a \omterm{def:m3:blockchains-for-traders:chain}{blockchain}'s state rather than by a company. An
\emph{on-chain order book}\index{on-chain order book} is a central limit order book whose orders, cancels
and matches are transactions executed and recorded by the chain's consensus.
\end{definition}

A market maker on a centralised venue sends thousands of messages a second and cancels most of its orders;
a message costs nothing and is processed in microseconds. On a general-purpose chain every message is a
transaction that pays gas (\cref{ch:m3:blockchains-for-traders}) and waits for a block, seconds away: quoting
costs money, cancels are slow, and anyone watching the \omterm{def:m3:blockchains-for-traders:mempool}{mempool} sees the cancel coming. Three consequences
follow. The chain must be built for the purpose, with blocks in fractions of a second and no per-message gas
for trading actions. The ordering of actions inside a block must be specified by the protocol, since the
block's proposer could otherwise sell the right to go first. And every function of the exchange that a
company would perform (the index, the \omterm{def:m3:perpetual-futures:index}{mark price}, liquidation, the backstop) must be performed by the
protocol and its \omterm{def:m3:blockchains-for-traders:pow}{validators}, in public.

\begin{definition}[App-chain]\label{def:m3:on-chain-order-books-and-perpetual-dexs:appchain}
An \emph{app-chain}\index{app-chain} is a \omterm{def:m3:blockchains-for-traders:chain}{blockchain} built to run one application, here an exchange, with its
own \omterm{def:m3:blockchains-for-traders:pow}{validators}, consensus parameters and transaction rules chosen for that application rather than for
general computation.
\end{definition}

\section{Hyperliquid in detail}

Hyperliquid is an \omterm{def:m3:on-chain-order-books-and-perpetual-dexs:appchain}{app-chain} whose state includes the order books, margin and positions of its perpetual and
spot markets (it calls this part HyperCore), alongside a general-purpose smart-contract environment.

\begin{dated}{2026-09}{Consensus, latency and throughput}\label{dat:m3:on-chain-order-books-and-perpetual-dexs:hl}
Hyperliquid's documentation describes a proof-of-stake chain secured by HyperBFT, a variant of HotStuff consensus,
with the order books and margin state held in the chain's state rather than off chain. It reports a median
end-to-end latency of 0.2 seconds and a 99th percentile of 0.9 seconds for an order from a co-located client, and
a throughput of about 200\,000 orders a second, limited by execution. It states that there is no designated
market-maker programme, and no special rebates, fees or latency advantages.
\end{dated}

\begin{definition}[In-block priority rule]\label{def:m3:on-chain-order-books-and-perpetual-dexs:priority}
An \emph{in-block priority rule}\index{in-block priority rule} is the protocol rule that sorts the trading actions
of a block before they are applied to the order books, overriding the order in which the block's proposer received
or listed them.
\end{definition}

\begin{dated}{2026-09}{The in-block ordering of one venue}\label{dat:m3:on-chain-order-books-and-perpetual-dexs:order}
Hyperliquid's documentation states that within a block, actions are sorted in each consensus batch (usually one or
two per block) as: first, actions that send no GTC or IOC order to any book; second, cancels; third, actions that send at
least one GTC or IOC order. Within each category actions keep the order proposed by the block's proposer. Its books match
in price-time priority, with prices on a tick grid and sizes on a lot grid; \omterm{def:m3:centralised-exchanges:postonly}{post-only orders} (ALO) rest without executing.
\end{dated}

\begin{omfigure}
\centering
\begin{tikzpicture}[font=\small,
  c/.style={draw,thick,rounded corners=2pt,align=center,minimum height=0.9cm,text width=3.3cm},
  arr/.style={-{Stealth[length=2.2mm]},thick}]
\node[draw,thick,rounded corners=2pt,align=left,text width=4.3cm,fill=gray!6] (in) at (-0.4,0) {\footnotesize actions as proposed:\\
IOC buy (taker)\\cancel stale ask (maker)\\post-only new ask (maker)\\GTC sell (trader)};
\node[c,draw=omDef,fill=omDef!8] (c1) at (5.6,1.3) {1.\ no GTC/IOC:\\post-only quotes};
\node[c,draw=omThm,fill=omThm!8] (c2) at (5.6,0) {2.\ cancels};
\node[c,draw=omExo,fill=omExo!8] (c3) at (5.6,-1.3) {3.\ GTC and IOC orders};
\node[c,draw=omProp,fill=omProp!8] (bk) at (10.6,0) {order book,\\price-time priority};
\draw[arr] (in.east) -- (c1.west); \draw[arr] (in.east) -- (c2.west); \draw[arr] (in.east) -- (c3.west);
\draw[arr] (c1.east) -- (bk.west); \draw[arr] (c2.east) -- (bk.west); \draw[arr] (c3.east) -- (bk.west);
\end{tikzpicture}

\omcaption{An \omterm{def:m3:on-chain-order-books-and-perpetual-dexs:priority}{in-block priority rule} of the kind \cref{dat:m3:on-chain-order-books-and-perpetual-dexs:order} describes: the
block's actions are sorted into three classes and applied class by class, in proposed order within each class. A maker's
cancel and new post-only quote are applied before any taker's order in the same block. Schematic.}\label{fig:m3:on-chain-order-books-and-perpetual-dexs:classes}
\end{omfigure}

The rule is a deliberate transfer of value from takers to makers. On a centralised venue a taker who reacts faster than a
maker's cancel picks off the stale quote; here, if both reach the same block, the cancel wins, as does the maker's new
post-only quote. Makers can quote tighter because they are picked off less, and a taker who wants to hit a stale quote must
get into an earlier block than the maker's cancel. The tutorial replays a thousand blocks: with the taker arriving first half
the time, arrival order lets it pick off 610 lots, 940 ticks of the maker's value; the class rule, none.

\begin{omfigure}
\begin{tikzpicture}
\begin{axis}[width=11cm,height=5cm,
  xlabel={block},ylabel={maker's cumulative loss, ticks},
  tick label style={font=\small},label style={font=\small},
  xmin=0,xmax=1000,ymin=0,ymax=1000,grid=major,grid style={gray!25},
  legend columns=2,legend style={font=\footnotesize,at={(0.5,-0.3)},anchor=north,draw=none,fill=none,/tikz/every even column/.append style={column sep=8pt}},
  table/col sep=comma]
\addplot[omThm,very thick] table[x=block,y=arrival]{figdata/markets-3/21-on-chain-order-books-and-perpetual-dexs/pickoff.csv};\addlegendentry{arrival order}
\addplot[omDef,very thick,dashed] table[x=block,y=cancels_first]{figdata/markets-3/21-on-chain-order-books-and-perpetual-dexs/pickoff.csv};\addlegendentry{quotes and cancels before takers}
\end{axis}
\end{tikzpicture}

\omcaption{A maker's losses to stale-quote pick-offs over a thousand blocks under two in-block rules, with fair value moving a
normal two ticks a block and the taker's order arriving before the maker's cancel half the time. Under the class rule the
maker's cancel and requote always come first. Synthetic replay. Data: the chapter's tutorial.}\label{fig:m3:on-chain-order-books-and-perpetual-dexs:pickoff}
\end{omfigure}

The rest of the exchange is on chain too. Prices come from \omterm{def:m3:blockchains-for-traders:pow}{validators} rather than from the venue's own trades alone, so that the
thin books of a young venue cannot be used to move its own references.

\begin{dated}{2026-09}{Oracle, mark and liquidation at one venue}\label{dat:m3:on-chain-order-books-and-perpetual-dexs:oracle}
Hyperliquid's \omterm{def:m3:blockchains-for-traders:pow}{validators} each publish, about every three seconds, a spot \omterm{def:m3:blockchains-for-traders:contract}{oracle} price per perpetual: a weighted median of spot mid
prices on Binance, OKX, Bybit, Kraken, KuCoin, Gate.io, MEXC and Hyperliquid, with weights 3, 2, 2, 1, 1, 1, 1 and 1 (external sources
only for assets whose main spot liquidity is elsewhere); the clearinghouse uses the stake-weighted median of the \omterm{def:m3:blockchains-for-traders:pow}{validators}' prices.
The \omterm{def:m3:perpetual-futures:index}{mark price} is the median of the \omterm{def:m3:blockchains-for-traders:contract}{oracle} price plus a 150-second moving average of the basis, the median of the venue's best bid,
best ask and last trade, and a weighted median of perpetual mids on five external venues. Maintenance margin is half the initial
margin at maximum leverage (1.25\% to 16.7\%); liquidations are first sent to the book as market orders (20\% at a time for positions
above 100\,000 USDC); below two thirds of maintenance margin a position is taken over by the liquidator vault, a strategy of the
protocol's vault, HLP, whose depositors share its P\&L.
\end{dated}

\begin{definition}[Liquidity vault]\label{def:m3:on-chain-order-books-and-perpetual-dexs:vault}
A \emph{liquidity vault}\index{liquidity vault} is a pool of depositors' capital run by a protocol, or by a manager under protocol rules,
that makes markets, takes over liquidated positions or supplies capital to a venue, sharing its P\&L among its depositors.
\end{definition}

A vault that backstops liquidations is the venue's \omterm{def:m3:margin-liquidation-and-loss-allocation:fund}{insurance fund} and its market maker of last resort, owned by whoever deposits in it.
Its depositors earn the liquidation buffer on average and bear the tail when a position it inherits cannot be closed at a sensible price.

\section{Incidents}

Stress tests the design where a centralised venue's staff would step in. In the 10 October 2025 sell-off
(\cref{ch:m3:margin-liquidation-and-loss-allocation}), CoinDesk Research measured the largest contraction of open interest on Hyperliquid,
from USD~14 billion to 6 billion, a fall of 57\%. The venue's \omterm{def:m3:margin-liquidation-and-loss-allocation:adl}{auto-deleveraging} closed profitable positions of other traders, and an analysis
by Tarun Chitra of Gauntlet, reported by Protos, argued that its queue-based ranking closed about USD~650 million more of them than was
necessary. Earlier, on 12 March 2025, a trader holding a 50 times leveraged long of 113\,000 ether withdrew margin until the position was
liquidated into the vault: the trader kept about USD~1.8 million of profit and the vault lost about 4 million; the venue then cut its
maximum leverage to 40 times on bitcoin and 25 on ether.

A second kind of incident is specific to a venue whose backstop is a public vault. A trader opens a large position in a thin market, lets
it be liquidated so that the vault inherits it, and then moves the external spot prices that feed the \omterm{def:m3:blockchains-for-traders:contract}{oracle}, on smaller venues where that is
cheap; the vault, holding a position it cannot close without moving the book, marks a growing loss. The protocol's defences are an
\omterm{def:m3:on-chain-order-books-and-perpetual-dexs:oicap}{open-interest cap} per asset and per account, margin tiers that raise maintenance margin with size, and, in the last resort, governance: the
\omterm{def:m3:blockchains-for-traders:pow}{validators} can vote to delist a market.

\begin{definition}[Open-interest cap]\label{def:m3:on-chain-order-books-and-perpetual-dexs:oicap}
An \emph{open-interest cap}\index{open-interest cap} is a limit on the total open interest of a market, or on the position of one account in it,
beyond which the venue rejects orders that would increase it.
\end{definition}

\begin{proposition}[The loss of a vault that inherits a squeezed position]\label{prop:m3:on-chain-order-books-and-perpetual-dexs:squeeze}
A vault takes over a short position of notional $N$ with a fraction $b$ of it left as margin, and the price then rises by a fraction $s$ before the
vault can close. Its loss is $N(s - b)^+$. To keep the loss within $L$ for squeezes up to $s_{\max}$, the positions it may inherit must be capped at
$N \le L/(s_{\max} - b)$.
\end{proposition}

\begin{proof}
The short loses $Ns$ and the inherited margin covers $Nb$ of it; the bound follows by solving $N(s_{\max} - b) \le L$.
\end{proof}

\begin{omfigure}
\begin{tikzpicture}
\begin{axis}[width=11cm,height=5cm,
  xlabel={squeeze after takeover, \%},ylabel={vault loss, USD million},
  tick label style={font=\small},label style={font=\small},
  xmin=0,xmax=500,ymin=0,ymax=50,grid=major,grid style={gray!25},
  legend columns=3,legend style={font=\footnotesize,at={(0.5,-0.3)},anchor=north,draw=none,fill=none,/tikz/every even column/.append style={column sep=8pt}},
  table/col sep=comma]
\addplot[omThm,very thick] table[x=squeeze_pct,y=nocap]{figdata/markets-3/21-on-chain-order-books-and-perpetual-dexs/squeeze.csv};\addlegendentry{USD~10 million inherited}
\addplot[omDef,very thick,dashed] table[x=squeeze_pct,y=cap2]{figdata/markets-3/21-on-chain-order-books-and-perpetual-dexs/squeeze.csv};\addlegendentry{cap USD~2 million}
\addplot[omExo,very thick,dotted] table[x=squeeze_pct,y=cap1]{figdata/markets-3/21-on-chain-order-books-and-perpetual-dexs/squeeze.csv};\addlegendentry{cap USD~1 million}
\end{axis}
\end{tikzpicture}

\omcaption{Loss of a vault that inherits a short in a market whose maximum leverage is 3 times (maintenance margin 16.7\%, takeover at two thirds of it,
leaving 11.1\% of notional as margin), as a function of the rise that follows (\cref{prop:m3:on-chain-order-books-and-perpetual-dexs:squeeze}). Illustrative
positions. Data: the chapter's tutorial.}\label{fig:m3:on-chain-order-books-and-perpetual-dexs:squeeze}
\end{omfigure}

\begin{dated}{2026-09}{Delisting by validators}\label{dat:m3:on-chain-order-books-and-perpetual-dexs:delist}
Hyperliquid's documentation states that \omterm{def:m3:blockchains-for-traders:pow}{validators} vote on whether to delist validator-operated perpetuals; if they do, the perpetuals settle to the
one-hour time-weighted spot \omterm{def:m3:blockchains-for-traders:contract}{oracle} price before the scheduled voting time, all positions are settled and open orders cancelled.
The pattern above happened on 26 March 2025: a JELLY short passed to the vault, HLP, through liquidation, spot buying elsewhere took the vault's
unrealised loss to USD~13.5 million, and the \omterm{def:m3:blockchains-for-traders:pow}{validators} voted to delist the market, which settled at 0.0095 while about 0.50 was being fed to the
\omterm{def:m3:blockchains-for-traders:contract}{oracle}.
\end{dated}

A vote that settles a market at a price is the decentralised equivalent of an exchange's emergency powers. It can protect a vault from a squeeze; it
also shows that the market's rules can be changed by the \omterm{def:m3:blockchains-for-traders:pow}{validators} while positions are open, a governance risk that a trader prices in the same way
as a centralised venue's discretion.

\section{App-chain and oracle-priced designs}

Two other designs trade some of the on-chain book's properties for performance or simplicity.

\begin{dated}{2026-09}{An app-chain with an off-chain book}\label{dat:m3:on-chain-order-books-and-perpetual-dexs:dydx}
dYdX's documentation describes its chain (``v4'') as a proof-of-stake chain built on CometBFT and the Cosmos SDK, whose \omterm{def:m3:blockchains-for-traders:pow}{validators} keep the order book
in memory, off chain and outside consensus; when orders match, the block's proposer adds the matches to its proposed block, and a block is committed when
two thirds of the stake approve it. Proposers take turns in a stake-weighted round robin.
\end{dated}

\begin{definition}[Oracle-priced exchange]\label{def:m3:on-chain-order-books-and-perpetual-dexs:oracle}
An \emph{oracle-priced exchange}\index{oracle-priced exchange} is a perpetual or spot DEX that has no order book: traders open and close positions
against a \omterm{def:m3:automated-market-makers:amm}{liquidity pool} at a price supplied by an \omterm{def:m3:blockchains-for-traders:contract}{oracle}, with fees and price-impact charges set by the protocol, and the pool's depositors take the
other side of every trade.
\end{definition}

\begin{dated}{2026-09}{An oracle-priced perpetual exchange}\label{dat:m3:on-chain-order-books-and-perpetual-dexs:gmx}
GMX's documentation describes a decentralised spot and perpetual exchange on Arbitrum, Avalanche and MegaETH, with leverage up to 100 times, that routes
every order against its \omterm{def:m3:automated-market-makers:amm}{liquidity pools} and quotes the \omterm{def:m3:blockchains-for-traders:contract}{oracle} \omterm{def:m3:perpetual-futures:index}{index price} rather than relying on an order book or external market makers; it uses Chainlink
Data Streams \omterm{def:m3:blockchains-for-traders:contract}{oracles}, and liquidity providers earn 63\% of the fees from trading, liquidations, borrowing and swaps on Arbitrum and Avalanche.
\end{dated}

\begin{omfigure}
\centering
\begin{tikzpicture}[font=\footnotesize,
  h/.style={draw,thick,rounded corners=2pt,align=center,text width=2.9cm,minimum height=0.8cm},
  r/.style={align=center,text width=2.9cm}]
\node[h,draw=omDef,fill=omDef!8] at (0,0) {centralised\\exchange};
\node[h,draw=omThm,fill=omThm!8] at (3.5,0) {on-chain\\order book};
\node[h,draw=omProp,fill=omProp!8] at (7,0) {off-chain book,\\on-chain matches};
\node[h,draw=omExo,fill=omExo!8] at (10.5,0) {oracle-priced\\pool};
\node[r,anchor=north] at (0,-0.6) {book: company\\order: its clock\\custody: company};
\node[r,anchor=north] at (3.5,-0.6) {book: chain\\order: rule\\custody: chain};
\node[r,anchor=north] at (7,-0.6) {book: validators\\order: proposer\\custody: chain};
\node[r,anchor=north] at (10.5,-0.6) {book: none\\price: oracle\\custody: chain};
\end{tikzpicture}

\omcaption{Four designs of a derivatives venue, by who holds the book, who decides the order of actions, and who holds the collateral. Schematic.}\label{fig:m3:on-chain-order-books-and-perpetual-dexs:designs}
\end{omfigure}

The oracle-priced design removes order-book games and replaces them with \omterm{def:m3:blockchains-for-traders:contract}{oracle} games: the pool loses whenever a trader knows the price better than the
\omterm{def:m3:blockchains-for-traders:contract}{oracle}, as an \omterm{def:m3:automated-market-makers:amm}{automated market maker} loses to arbitrageurs (\cref{ch:m3:automated-market-makers}), and it must defend itself with fees, delays and limits on
open interest. The off-chain book restores speed but hands the ordering of matches to each block's proposer. For a trading firm the questions are those of
any venue (\cref{ch:m3:centralised-exchanges}), with different answers: who orders its messages, what happens to its collateral if the venue fails, and who
can change the rules while it is positioned.

\section{Tutorial: two orderings of one block}\label{tut:m3:on-chain-order-books-and-perpetual-dexs:replay}

\begin{tutorial}
\textbf{Goal.} Replay blocks of maker quotes, cancels and taker orders through a price-time-priority book under arrival order and under the class rule, and
count the stale quotes picked off. \textbf{End state:} \cref{fig:m3:on-chain-order-books-and-perpetual-dexs:pickoff,fig:m3:on-chain-order-books-and-perpetual-dexs:squeeze}.
\begin{tutsteps}
\item \textbf{The rule.} Sort a block's actions by class, keeping arrival order within each class.
\omcode{code/firm/blockbook/firm_blockbook.py}{83}{96}{The in-block ordering and the block runner.}\label{lst:m3:on-chain-order-books-and-perpetual-dexs:order}
\item \textbf{The replay.} A maker requotes every block; a taker attacks quotes that fair value has passed.
\omcode{code/markets-3/21-on-chain-order-books-and-perpetual-dexs/python/m3_perpdex.py}{12}{41}{The pick-off replay.}\label{lst:m3:on-chain-order-books-and-perpetual-dexs:replay}
\item \textbf{Run} \texttt{replay} with the rules \texttt{arrival} and \texttt{cancels\_first}, then \texttt{fig\_perpdex.py}.
\end{tutsteps}
\textbf{What to change next.} Let the taker arrive first 90\% of the time; let the maker's cancel fall in the next block with some probability and see what the class
rule still protects; add a second maker who never cancels.
\end{tutorial}

\section{Build: the block-batched book}\label{bld:m3:on-chain-order-books-and-perpetual-dexs:blockbook}

\begin{build}
\bfield{Purpose} To trade on block-based venues the miniature firm must predict how its actions will be ordered and matched; to backtest there it must replay
blocks exactly; and its risk must include the venue's \omterm{def:m3:on-chain-order-books-and-perpetual-dexs:oicap}{open-interest caps} and its \omterm{def:m3:perpetual-futures:index}{mark price}.
\bfield{Interface} \texttt{Action(kind, trader, side, price, qty, oid)} with kinds \texttt{alo}, \texttt{gtc}, \texttt{ioc}, \texttt{cancel};
\texttt{Book(oi\_cap)} with \texttt{apply(action)} and \texttt{best()}; \texttt{order\_block(actions, rule)} for \texttt{arrival} and
\texttt{cancels\_first}; \texttt{run\_block(book, actions, rule)}; \texttt{mark\_price(oracle\_basis, book\_median, external)}, the median of the three. Uses \texttt{firm.perp} and
\texttt{firm.liquidation} for margin.
\bfield{Rules} Integer ticks and lots; stable sorting, so that replays are deterministic; \omterm{def:m3:centralised-exchanges:postonly}{post-only orders} that would cross are rejected; orders that would take a
position beyond the cap are rejected.
\bfield{Acceptance tests} \texttt{code/firm/blockbook/tests/}: price-time priority and post-only rejection; a stale quote picked off under arrival order and protected under
the class rule; the \omterm{def:m3:on-chain-order-books-and-perpetual-dexs:oicap}{open-interest cap}; the median mark.
\bfield{Stretch} Consensus batches within a block; margin checks at order entry and at the resting side's fill; liquidation orders entering the block as IOCs.
\end{build}

\begin{omsources}
\item Hyperliquid documentation: HyperCore overview, order book, oracle, robust price indices, liquidations, protocol vaults, delisting, market making.
Accessed September 2026.
\item dYdX documentation, ``Intro to dYdX Chain Architecture''; GMX documentation, introduction. Accessed September 2026.
\item CoinDesk Research, 17 October 2025; Protos, 10 December 2025 (see \cref{ch:m3:margin-liquidation-and-loss-allocation}).
\item CoinDesk, ``Hyperliquid Loses \$4M After Whale's Over \$200M Ether Trade Unwinds'', 12 March 2025; ``HyperLiquid Delists JELLY
After Vault Squeezed in \$13M Tussle'', 26 March 2025.
\end{omsources}

\section{Exercises}

\begin{exercise}[$\star$]\label{exo:m3:on-chain-order-books-and-perpetual-dexs:1}
In one block arrive, in this order: an IOC buy at 101, a cancel of the ask at 101, and a post-only ask at 103. What executes under arrival order, and under the
class rule?
\end{exercise}

\begin{exercise}[$\star$]\label{exo:m3:on-chain-order-books-and-perpetual-dexs:2}
Why can a general-purpose chain with 12-second blocks and per-transaction gas not host a competitive order book?
\end{exercise}

\begin{exercise}[$\star$]\label{exo:m3:on-chain-order-books-and-perpetual-dexs:3}
What does the stake-weighted median of \omterm{def:m3:blockchains-for-traders:pow}{validators}' \omterm{def:m3:blockchains-for-traders:contract}{oracle} prices protect against?
\end{exercise}

\begin{exercise}[$\star\star$]\label{exo:m3:on-chain-order-books-and-perpetual-dexs:4}
An asset has a maximum leverage of 20 times. What is its maintenance margin under the rule of \cref{dat:m3:on-chain-order-books-and-perpetual-dexs:oracle}, and at
what equity is a position taken over by the vault?
\end{exercise}

\begin{exercise}[$\star\star$]\label{exo:m3:on-chain-order-books-and-perpetual-dexs:5}
A vault inherits a USD~4 million short at 3 times maximum leverage and the price then doubles. What does it lose?
\end{exercise}

\begin{exercise}[$\star\star$]\label{exo:m3:on-chain-order-books-and-perpetual-dexs:6}
Compare the risks a liquidity provider bears in an oracle-priced pool and in an order-book venue's vault.
\end{exercise}

\begin{exercise}[$\star\star\star$]\label{exo:m3:on-chain-order-books-and-perpetual-dexs:7}
\emph{Coding.} With \texttt{replay}, raise the probability that the taker arrives first to 0.9. How do the two rules compare?
\end{exercise}

\begin{exercise}[$\star\star\star$]\label{exo:m3:on-chain-order-books-and-perpetual-dexs:8}
\emph{Find the flaw.} ``The venue is decentralised, so no one can change a market's rules while I hold a position.''
\end{exercise}

\section{Problem: The Squeezed Vault}

\begin{problem}[{Weekend problem --- a thin market, a large short, a public backstop}]\label{pb:m3:on-chain-order-books-and-perpetual-dexs:1}
A thin \omterm{def:m3:blockchains-for-traders:key}{token} trades as a perpetual with a maximum leverage of 3 times. A trader opens a USD~10 million short and lets it be liquidated; the vault takes it over at two
thirds of the maintenance margin. The \omterm{def:m3:blockchains-for-traders:key}{token}'s spot price, on small external venues, is then pushed up.

\medskip\noindent\textbf{Part I --- The takeover.}
\begin{enumerate}
\item What is the maintenance margin, and what margin comes with the position when the vault takes it over?
\item Why could the position not be closed in the book before the takeover?
\item What does the vault lose if the price rises 100\% before it can close? 400\%?
\item Why can the trader profit although his short was liquidated?
\item What does the \omterm{def:m3:blockchains-for-traders:contract}{oracle}'s construction make the attacker do?
\end{enumerate}

\medskip\noindent\textbf{Part II --- The cap.}
\begin{enumerate}[resume]
\item What cap on a single position keeps the vault's loss within USD~5 million for squeezes up to 500\%?
\item What would a cap of USD~2 million have limited the loss to at 400\%?
\item Why should the cap depend on the depth of the external spot markets?
\item What other parameters could the venue tighten?
\item What does the vault's depositor need to know about these parameters?
\end{enumerate}

\medskip\noindent\textbf{Part III --- Governance.}
\begin{enumerate}[resume]
\item What does a \omterm{def:m3:blockchains-for-traders:pow}{validators}' vote to delist do to open positions?
\item Who gains and who loses from settling at the \omterm{def:m3:blockchains-for-traders:contract}{oracle}'s one-hour average?
\item Is the vote a failure of decentralisation or a feature of it?
\item How would a centralised venue have handled the same attack?
\item How does a trading firm price this governance risk?
\end{enumerate}

\medskip\noindent\textbf{Part IV --- Judgement.}
\begin{enumerate}[resume]
\item Should a vault backstop thin markets at all?
\item What makes an on-chain venue safer than a centralised one, and what does not?
\item Where would you want the firm's market making: on the order-book DEX, the off-chain book or the oracle-priced pool?
\item State the \emph{named result}: the vault's loss as a function of the squeeze for a USD~10 million position, and the cap that bounds it to USD~5 million at 500\%.
\item In one sentence: what is a public backstop?
\end{enumerate}
\end{problem}

\section{Interview questions}

\begin{interviewq}[$\star$ \iqroles{trader}]\label{iq:m3:on-chain-order-books-and-perpetual-dexs:1}
How does an \omterm{def:m3:on-chain-order-books-and-perpetual-dexs:priority}{in-block priority rule} that applies cancels before orders change a market maker's economics?
\end{interviewq}

\begin{interviewq}[$\star$ \iqroles{developer}]\label{iq:m3:on-chain-order-books-and-perpetual-dexs:2}
What changes in your order-management system when the venue is a \omterm{def:m3:blockchains-for-traders:chain}{blockchain} with sub-second blocks?
\end{interviewq}

\begin{interviewq}[$\star\star$ \iqroles{researcher}]\label{iq:m3:on-chain-order-books-and-perpetual-dexs:3}
Compare a weighted median of external prices with a volume-weighted average as an \omterm{def:m3:blockchains-for-traders:contract}{oracle}.
\end{interviewq}

\begin{interviewq}[$\star\star$ \iqroles{risk}]\label{iq:m3:on-chain-order-books-and-perpetual-dexs:4}
What limits would you set for a firm's exposure to a \omterm{def:m3:on-chain-order-books-and-perpetual-dexs:perpdex}{perpetual DEX}'s vault or pool?
\end{interviewq}

\begin{interviewq}[$\star\star$ \iqroles{trader, researcher}]\label{iq:m3:on-chain-order-books-and-perpetual-dexs:5}
How would you attack an \omterm{def:m3:on-chain-order-books-and-perpetual-dexs:oracle}{oracle-priced exchange}, and how would you defend it?
\end{interviewq}

\begin{interviewq}[$\star\star\star$ \iqroles{developer}]\label{iq:m3:on-chain-order-books-and-perpetual-dexs:6}
Design a deterministic replay system for a block-based venue's order books from its published blocks.
\end{interviewq}
